\chapter{Минимальный носитель спина три и препятствие прямоугольного блока}

Предыдущая глава доказала, что текущая конечная операторная алгебра не
содержит представления спина три. Внешний симметрический куб триплета
содержит его ровно один раз, но сам по себе ещё не задаёт конечную цепь.
Здесь проверяется минимальное встраивание, использующее только два уже
известных семейных типа: триплет $V_1$ и пятёрку
$V_2=\operatorname{Sym}^2_0(V_1)$ параметра порядка $Q$.

\section{Литературная и проектная предпосылка}

Разложение тензорных произведений $SO(3)$ задаётся стандартными
коэффициентами Клебша--Гордана; современное систематическое изложение
приведено, например, в \path{arXiv:1910.06821}. Для спинов один и два
имеем
\begin{equation}
 V_2\otimes V_1=V_1\oplus V_2\oplus V_3.
 \label{eq:v6-spin-one-two-clebsch-gordan}
\end{equation}
Представление $V_3$ является семимерным пространством полностью
симметричных бесследовых тензоров третьего ранга, используемым в явных
моделях $SO(3)\to A_4$; см. \path{arXiv:0910.4392} и
\path{arXiv:2607.12366}.

Конечные цепи и их внедиагональные операторы являются стандартной частью
диаграмм Краевского и сетей спектральных троек; см.
\path{arXiv:hep-th/9701081} и \path{arXiv:1301.3480}. Эти источники
разрешают архитектуру, но не определяют, какой угол кривизны должен войти
в действие настоящего проекта.

\section{Каноническая семёрка внутри отображений}

Рассмотрим вещественное пространство
\begin{equation}
 \operatorname{Hom}(V_1,V_2),
 \qquad \dim=5\cdot3=15.
\end{equation}
Пусть $C$ --- квадратичный казимир совместного действия $SO(3)$ на этом
пространстве. Его спектр равен
\begin{equation}
 \Spec C=\{2^{(3)},6^{(5)},12^{(7)}\},
 \label{eq:v6-hom-casimir-spectrum}
\end{equation}
что соответствует спинам один, два и три. Поэтому проектор на искомую
семёрку задаётся без свободного коэффициента:
\begin{equation}
 P_{(3)}=\frac{(C-2I)(C-6I)}{(12-2)(12-6)}
 =\frac{(C-2I)(C-6I)}{60}.
 \label{eq:v6-spin-three-casimir-projector}
\end{equation}
Машинная проверка даёт
\begin{equation}
 \rank P_{(3)}=7,
 \qquad
 \|P_{(3)}^2-P_{(3)}\|=1.6\cdot10^{-15}.
\end{equation}

Для $T_{ijk}\in\operatorname{Sym}^3_0(V_1)$ определим отображение
$Z_T:V_1\to V_2$ формулой
\begin{equation}
 (Z_T e_k)_{ij}=T_{ijk}.
 \label{eq:v6-tensor-to-triplet-quintet-arrow}
\end{equation}
Поскольку срез $T_{ij k}$ симметричен и бесследов по $i,j$, он
действительно лежит в $V_2$. Это отображение является изометрией
\begin{equation}
 \operatorname{Sym}^3_0(V_1)
 \xrightarrow{\ \simeq\ }
 \operatorname{im}P_{(3)}
 \subset\operatorname{Hom}(V_1,V_2).
 \label{eq:v6-spin-three-arrow-isometry}
\end{equation}
Остаток изометрии равен $8.8\cdot10^{-16}$, а остаток проекции на
семёрку --- $1.6\cdot10^{-15}$.

\begin{theorem}[Минимальный связанный носитель]
Искомое поле спина три канонически и с кратностью один реализуется как
семимерное подрасслоение
$\operatorname{im}P_{(3)}\subset\operatorname{Hom}(V_1,V_2)$. Оно
использует одну уже существующую семейную связность и не требует нового
калибровочного множителя.
\end{theorem}

\section{Тетраэдрическая кривизна как квадрат стрелки}

Сопряжённый квадрат на триплетном углу равен
\begin{equation}
 (Z_T^\dagger Z_T)_{kl}
 =\sum_{ij}T_{ijk}T_{ijl}
 =T_{kij}T_{lij}=A_{kl}(T).
 \label{eq:v6-arrow-square-tensor-contraction}
\end{equation}
Следовательно,
\begin{equation}
 \mathcal F_{1}(T)
 =Z_T^\dagger Z_T-\frac{v_T^2}{3}I_3
 =\mu_T(T).
 \label{eq:v6-triplet-corner-tetrahedral-curvature}
\end{equation}
Это не аналогия: остаток тождества в вычислении равен нулю.

На каноническом тетраэдрическом вакууме
$v_T^2=32/9$ три сингулярных числа $Z_{T_\star}$ совпадают и равны
\begin{equation}
 \sqrt{\frac{32}{27}}=1.0886621079\ldots,
\end{equation}
поэтому кривизна~\eqref{eq:v6-triplet-corner-tetrahedral-curvature}
обращается в нуль.

Гессиан нормированного квадрата
\begin{equation}
 \frac13\Tr\mathcal F_1(T)^2
\end{equation}
на семимерном подпространстве имеет тот же спектр, что и независимый
тензорный расчёт:
\begin{equation}
 \left\{0^{(3)},
 \left(\frac{512}{405}\right)^{(3)},
 \frac{256}{81}\right\}.
 \label{eq:v6-embedded-spin-three-hessian}
\end{equation}
Массовая матрица одной семейной связности снова равна
$(128/9)I_3$.

Кроме того, усиление физическим пакетом $H_{15}$ не меняет
нормировку:
\begin{equation}
 \frac1{45}\Tr_{V_1\otimes H_{15}}
 (\mathcal F_1^2\otimes I_{15})
 =\frac13\Tr_{V_1}\mathcal F_1^2.
 \label{eq:v6-spin-three-H45-trace-restriction}
\end{equation}
Максимальный численный остаток равен $1.8\cdot10^{-15}$.

\begin{proposition}[Проход триплетного угла]
Каноническая стрелка $V_1\to V_2$, ограниченная проектором
$P_{(3)}$, точно воспроизводит тетраэдрическую матричную кривизну,
правильный гессиан, массу связности и нормированный общий след.
\end{proposition}

\section{Препятствие буквального самосопряжённого блока}

Нельзя, однако, автоматически объявить $Z_T$ новым конечным оператором.
Буквальный нечётный самосопряжённый блок имеет вид
\begin{equation}
 D_T=
 \begin{pmatrix}
 0&Z_T^\dagger\\
 Z_T&0
 \end{pmatrix}
 \quad\text{на}\quad V_1\oplus V_2.
 \label{eq:v6-rectangular-spin-three-dirac-block}
\end{equation}
Он самосопряжён и антикоммутирует с градуировкой, но
$Z_T$ является матрицей $5\times3$. Поэтому
\begin{equation}
 \rank D_T=6,
 \qquad
 \dim\ker D_T=2.
 \label{eq:v6-rectangular-block-kernel}
\end{equation}
Два нуля лежат в дополнении образа $Z_T$ внутри пятёрки. После усиления
$H_{15}$ они превращаются в тридцать нулевых направлений.

Второй угол также не может иметь изотропную нулевую кривизну. Матрица
$Z_TZ_T^\dagger$ имеет спектр
\begin{equation}
 \left\{\frac{32}{27},\frac{32}{27},\frac{32}{27},0,0\right\}.
\end{equation}
Никакой скалярный сдвиг на $V_2$ не делает её пропорциональной $I_5$.
Даже оптимальный центральный сдвиг оставляет ненулевой остаток
$1.2983053215\ldots$.

\begin{nogocriterion}[Запрет полного двухуглового квадрата]
Обычный полный след самосопряжённого блока
\eqref{eq:v6-rectangular-spin-three-dirac-block} не сводится к одному
тетраэдрическому квадрату. Он неизбежно видит двухмерный коядровый сектор
пятёрки. Поэтому нельзя выдавать проход триплетного угла за готовое
встраивание в буквальный конечный оператор.
\end{nogocriterion}

\section{Новая точная граница}

Минимальный представительный вопрос решён положительно: семёрка не
произвольна и не требует второй связности. Но действие ещё не замкнуто.
Требуется вывести из прежней градуировки или условного ожидания именно
одностороннюю относительную кривизну
\begin{equation}
 \mathcal F_1=Z_T^\dagger Z_T-\frac{v_T^2}{3}I_3,
\end{equation}
не добавляя вручную проектор на удобный угол и не превращая коядро
$V_2$ в новые физические нулевые состояния.

Следующий гейт должен сопоставить этот угол с уже существующими в проекте
коммутирующим квадратом, отображением момента и градуированной высотой.
Если каноническая проекция не возникнет, буквальная конечная
спектральная ветвь спина три будет закрыта, хотя эффективная бозонная
полевая реализация сохранится.

\section{Вердикт}

\begin{protocol}
\begin{enumerate}
 \item разложение $V_2\otimes V_1=V_1\oplus V_2\oplus V_3$:
 \textbf{проверено};
 \item кратность семёрки: \textbf{один};
 \item канонический проектор Казимира: \textbf{получен};
 \item новое калибровочное поле: \textbf{не требуется};
 \item квадрат стрелки на триплетном углу: \textbf{$\mu_T$ точно};
 \item полный семимерный гессиан: \textbf{$3$ нуля, $4$ плюса};
 \item масса семейной связности: \textbf{положительна};
 \item нормировка на $H_{45}$: \textbf{совпадает};
 \item буквальный самосопряжённый блок: \textbf{имеет два нуля};
 \item полный двухугловой квадрат: \textbf{не проходит};
 \item односторонняя угловая проекция уже выведена: \textbf{нет};
 \item рождение материи: \textbf{не закрыто};
 \item следующий гейт: \textbf{родительский вывод угловой кривизны
 спина три}.
\end{enumerate}
\end{protocol}

Машинный сертификат сохранён в
\path{s2t_v6_bosonic_defect_minimal_spin_three_carrier_embedding_gate_results.json}.